Noter fra mødet 09-02:
- Fedt med dildo sugekop
- Wheatstone, 2 active og 2 inactive
- Piezoresistive actuator er en mulighed, start med primitiv testsetup
- Zigbee router yes
- Gør det kompakt men hvis batteriet skal om på den anden side er det fint
- Vi har 30.000 jyske dollars!

Noter fra mødet 16-02:
(Gabi var der ikke)
- Konstant modstandsfald var ikke et problem med conductive PLA. Prøv at undersøg det noget mere

Noter fra mødet 23-02:
(spontant møde med Gabi kl.10 og opfølgning med Anders kl.11)
- Konstant modstandsfald er material creep, ikke et stort problem når der er gået mere end en dag   og i kombi med wheatstone og måske moving average filter eller lign.
- RF, bare tag det nemme valg
- Power analyzer, Gabi finder noget
- Mekanisk setup:
    - Alu plade, prøv det.
    - Lav ny setup også til plastik
    - Det er ikke et problem at den hænger i fri luft når vi presretcher

Noter fra mødet 02-03:
(Møde kl.11 aflyst, erstattet af spontant møde med Gabi)
- Projektplan kan bare forklare indledningsvis at 3D printed sensor er den bedste løsning i         forhold til krav og der fokuseres dermed på det.

Noter fra mødet 09-03:
(intet møde, Gabi er syg, han finder os senere på ugen)
- Bare køb en power analyzer
- PCB review, der er lidt små ændringer men overordnet ser det fedt ud.
- Overvej coincell batteri igen

Noter fra mødet 16-03:
- Batteri og PCB tager vi senere, men ser fint ud
- Mounting skal ikke tage for meget fokus, kør den ind med tape
- General catch up
- 1.8 V kontra 3 V
    - 1.8 V kan bruges mere universelt
- Skal vi overhovedet pre-stretch? Hvad er argumentet, er det bare pga den gamle idé at vi tror     vi skal pre-stretch

Noter fra mødet 23-03:
(møde med Anders d.23)
- PCB og schematic har han ikke så meget at sige til men det ser fint ud
    - RF er fint ikke at gå for meget i dybden med, i forhold til Vestas krav burde 20 m være fint
    - Den cap på ADC'en, hvad gør den?
- Data ser pisse fedt ud. Umiddelbart fint ikke at pre-stretche den. Prøv at kør den ind med 2      Keithly, eller 6-7 stykker
(møde med Gabi d.24)
- PCB og schematic er super, prøv at ryk rundt på connecterne bl.a. ved at bruge den anden ADC      mulighed på MCU'en
- Det er super fint at købe impedance matching gennem JLBC
- Det er fedt at vi har fået systemets signaler til at snakke sammen, sensor ID må også gerne       komme på.
- Det er fint at have nogle test der siger at vi ikke skal pre-stretche og simplificere processen   på den måde.
- De lange test er fede og klima kammeret bliver spændende, vi bør også lave en så lang som mulig   cyclus test måske hen over påsken.

Noter fra mødet 30-03:
- Det påske bro

Noter fra mødet 06-04:
- Det stadig påske bro

Noter fra mødet 13-04:
(Gabi er syg men har givet go til PCB assembly, remote)
Anders skipper også mødet
Spontant møde med Gabi torsdagen efter:
- Fedt at PCB virker!!!
    - I kan også simulere om LDO er nødvendig hvis datasheet er gode
    - Tjek op på ADC'ens input impedans, den kan påvirke vores sensor måling som en ekstra          "hidden" voltage divider
- Påske testen, prøv at signal process det noget mere så vi måske kan se en relation. (cross        correlate)
- Klima test, giv den gas.
    - start i positive temp og kør cyclus, på et andet tidspunkt kør ekstrem -40 til 120 og se      hvad der sker med sensoren. 
    - Kør humidity også, men vi regner ikke med at se noget.
    - Lav wheatstone og kør den test. Hvis den er god, kør den med DC motoren også.

Noter fra mødet 20-04:
- Problemer med at flashe kode er pga clash med 3.3V og 1.8V, det kan fikses enten med software     eller ekstern hardware
- Klima kammeret trængte bare til service så vi bruger det bare, Gabi confirmed
- Prøv at print en sensor i TPU og i PLA (conductive) og test, ellers må vi bare sige at TPU er     besværligt og under andre conditions vil det måske blive bedre

Noter fra mødet 27-04:
(Gabi syg igen)
- Prøv at kør testene med samme sensor hvor at de starter på samme værdi
- Print på stål plade med tape ovenpå
- Health tech har en "instron" maskine
- Input fra student seminar
    - Skrue test setup med torquemeter (moment nøgle)

Noter fra mødet 04-05:
- Prøv at printe conductive PLA på metalplade og se drift
- Instron (Gabi siger glem det)
- Nyt pcb behøver vi ikke, bare et ny BOM
- Teste pcb komponenter, om vi behøver dem alle.
- Test zigbee kommunikation, RSSI i HA. Ved -80dBm begynder man måske at miste data

Noter fra mødet 11-05:
(Gabi er på ferie så kun Anders)
- Wheatstone-bridge bias skal ikke være lig med sensing rest, vi vil hellere have un-balanced da vi kun sensor i én retning. Og der skal bias være den relative midte mellem min og max ikke den lineære:
    - Ex, sensing range på 10 til 40. Lineært ville være 25, men relativt ville være 20 (10*2*2)
- Humidity test, den er påvirket under humidity påvirkning men ikke bagefter, og det er bare det.

Noter fra mødet 18-05:
- Lav acedemic plot med hysterese også
- Skriv strain i micro strain i stedet for
- Lav wheatstone plot også. mV/V.
- Køb fast silikone med afskærmning til antenne og connectors

Noter fra mødet 25-05:
- 

Noter fra mødet 01-06:
- 

Noter fra mødet 08-06:
- 

Noter fra mødet 15-06:
- 

Noter fra mødet 22-06:
- 

Noter fra mødet 29-06:
- 